\scanpage{62}
% archive foliation: 61
(au dessus duquel $X^*$ est étale \deleted{\unsure{galois}} principal de groupe $\underline V$), tout comme $E$ sur $E_0$) mais non au voisinage des pts de $D_J\subset\Sigma_J$ (vu que $E_0/\Sigma_J$ y est ramifié, mais non $E/X$ en les points de $\Delta=X\mid D_J$, or si c'était cartésien il en résulterait que $E_0/\Sigma_J$ est étale sur les pts de $D_J$ aussi, par descente plate\ldots). \unsure{Quand|Quant} $=$ il doit être facile de vérifier que le diagramme total se reconstitue à partir de $E_0$ ($\mapsto$ \inserted{$\underline J$ et} $\Sigma=\Sigma_{\underline J}$, et $E_0\longrightarrow\Sigma$, \inserted{et de section $t$ de $\Sigma^*$}) et du revêtement \deleted{principal} $X$ de $\Sigma$ avec groupe \inserted{\unsure{dég.}} $\underline V$: la donnée de $(E,\sigma)$ i.e. de $(E_0,I)$ \deleted{doit être} \inserted{est} équivalente aussi à $E_0$ et du revêtement $X$ (d'où en fait \deleted{$\underline I$}\inserted{$\underline K$} et $\underline I\simeq\underline K\wedge(X\mid t)\unsure{'|^{\circ}}\ldots$)

Cela montre que finalement $(E,\sigma)$ i.e. le $I$ est un peu de la poudre aux yeux, on va prendre pour le reste $E=E_0$, $\sigma=\sigma_0$, donc $I=\underline V={}_2E_0=M_0$, torseur trivial. On a donc \deleted{$\unsure{\Sigma|I}$} $X=\Sigma_J$ et on trouve donc
